% META: {"id": "1NSI-TAB-COURS-02", "chapitre": "1NSI-TABLES", "type_objet": "cours", "capacites": ["1NSI-TABLES-C2"], "mode_creation": "ex_nihilo", "sources_inspiration": [], "status": "approved", "capacites_codes": ["C2"]}
\section{Recherche dans une table}

\propriete[Recherche par critère]{
  Rechercher les lignes d'une table vérifiant un critère consiste à parcourir la table et à
  ne conserver que les lignes pour lesquelles une \textbf{expression booléenne} (le
  critère) est vraie. On peut combiner plusieurs conditions avec \lstinline{and},
  \lstinline{or}, \lstinline{not}.
}

\begin{codereference}{Recherche par critère (compréhension de liste)}
\begin{python}
adherents = [
    {"nom": "Ali", "age": 16, "categorie": "junior", "cotisation": 80},
    {"nom": "Sami", "age": 15, "categorie": "junior", "cotisation": 80},
    {"nom": "Nour", "age": 17, "categorie": "junior", "cotisation": 80},
    {"nom": "Yasmine", "age": 22, "categorie": "senior", "cotisation": 120},
    {"nom": "Karim", "age": 19, "categorie": "senior", "cotisation": 120},
]

def rechercher(table, critere):
    """Renvoie les lignes de `table` pour lesquelles critere(ligne) est vrai."""
    return [ligne for ligne in table if critere(ligne)]

juniors_16plus = rechercher(adherents, lambda ligne: ligne["categorie"] == "junior" and ligne["age"] >= 16)
print([ligne["nom"] for ligne in juniors_16plus])  # ['Ali', 'Nour']
\end{python}
\end{codereference}

% BEGIN-VERIFY
% adherents = [
%     {"nom": "Ali", "age": 16, "categorie": "junior", "cotisation": 80},
%     {"nom": "Sami", "age": 15, "categorie": "junior", "cotisation": 80},
%     {"nom": "Nour", "age": 17, "categorie": "junior", "cotisation": 80},
%     {"nom": "Yasmine", "age": 22, "categorie": "senior", "cotisation": 120},
%     {"nom": "Karim", "age": 19, "categorie": "senior", "cotisation": 120},
% ]
% def rechercher(table, critere):
%     return [ligne for ligne in table if critere(ligne)]
% juniors_16plus = rechercher(adherents, lambda ligne: ligne["categorie"] == "junior" and ligne["age"] >= 16)
% assert [ligne["nom"] for ligne in juniors_16plus] == ["Ali", "Nour"]
% END-VERIFY

\propriete[Recherche de doublons]{
  Une table contient un \textbf{doublon} lorsque plusieurs lignes ont la même valeur pour
  une colonne censée être unique (par exemple un identifiant). On détecte les doublons en
  comptant les occurrences de chaque valeur de cette colonne.
}

\begin{codereference}{Détecter les doublons d'une colonne}
\begin{python}
def doublons(table, colonne):
    """Renvoie les valeurs de `colonne` apparaissant plus d'une fois dans `table`."""
    valeurs = [ligne[colonne] for ligne in table]
    return [v for v in set(valeurs) if valeurs.count(v) > 1]

adherents_avec_doublon = adherents + [{"nom": "Ali", "age": 16, "categorie": "junior", "cotisation": 80}]
print(doublons(adherents_avec_doublon, "nom"))  # ['Ali']
print(doublons(adherents, "nom"))                # []
\end{python}
\end{codereference}

% BEGIN-VERIFY
% adherents = [
%     {"nom": "Ali", "age": 16, "categorie": "junior", "cotisation": 80},
%     {"nom": "Sami", "age": 15, "categorie": "junior", "cotisation": 80},
% ]
% def doublons(table, colonne):
%     valeurs = [ligne[colonne] for ligne in table]
%     return [v for v in set(valeurs) if valeurs.count(v) > 1]
% avec_doublon = adherents + [{"nom": "Ali", "age": 16, "categorie": "junior", "cotisation": 80}]
% assert doublons(avec_doublon, "nom") == ["Ali"]
% assert doublons(adherents, "nom") == []
% END-VERIFY

\propriete[Test de cohérence]{
  Un \textbf{test de cohérence} vérifie qu'une table respecte des contraintes attendues
  (types corrects, valeurs dans un intervalle plausible, absence de champ vide...). C'est
  une étape essentielle avant tout traitement automatisé de données réelles, souvent
  imparfaites.
}

\begin{codereference}{Test de cohérence simple}
\begin{python}
def table_coherente(table):
    """Verifie que chaque age est un entier positif raisonnable (0 a 120)."""
    for ligne in table:
        if not isinstance(ligne["age"], int) or not (0 <= ligne["age"] <= 120):
            return False
    return True

print(table_coherente(adherents))                                    # True
print(table_coherente(adherents + [{"nom": "X", "age": -5,
                                     "categorie": "junior", "cotisation": 80}]))  # False
\end{python}
\end{codereference}

% BEGIN-VERIFY
% adherents = [{"nom": "Ali", "age": 16, "categorie": "junior", "cotisation": 80}]
% def table_coherente(table):
%     for ligne in table:
%         if not isinstance(ligne["age"], int) or not (0 <= ligne["age"] <= 120):
%             return False
%     return True
% assert table_coherente(adherents) is True
% incoherente = adherents + [{"nom": "X", "age": -5, "categorie": "junior", "cotisation": 80}]
% assert table_coherente(incoherente) is False
% END-VERIFY

\erreurFrequente{Rechercher des doublons en comparant chaque ligne \textbf{entière}
(\lstinline{ligne1 == ligne2}) alors que la consigne porte sur une seule colonne (comme
\lstinline{nom}). Deux lignes peuvent partager le même nom sans être identiques sur les
autres colonnes (âge différent par exemple) : il faut comparer uniquement la colonne
concernée.}

\alamachine{Ajoute une ligne dupliquée (même \lstinline{nom} qu'une ligne existante, mais
\lstinline{age} différent) à la table \lstinline{adherents}. Utilise \lstinline{doublons}
pour la détecter, puis écris une fonction \lstinline{supprimer_doublons(table, colonne)}
qui ne garde que la première occurrence de chaque valeur de \lstinline{colonne}.}
